\documentclass[12pt]{article}
\usepackage[utf8]{inputenc}
\usepackage[T1]{fontenc}
\usepackage{hyperref}
\title{Miwa Language Resources for K-12 Curriculum Development}
\author{Compiled by Claude Sonnet 5 (claude-sonnet-5) for Andr\'as Kornai}
\date{}
\begin{document}
\maketitle
\section{Introduction}
Miwa (ISO 639-3 vmi, Glottocode miwa1242) is listed as extinct and belongs to the Worrorran family of the Kimberley region, Western Australia. Sources differ on whether it is a separate language, a dialect of Wunambal, or the same as Gonin. No heritage or revitalization materials specific to Miwa were found.
\section{Primary Educational Resources}
None found for Miwa.
\section{Linguistic and Reference Materials}
\begin{itemize}
\item Glottolog entry for Miwa: \url{https://glottolog.org/resource/languoid/id/miwa1242}
\item AIATSIS Austlang record: \url{https://collection.aiatsis.gov.au/austlang/language/K44}
\item Wikipedia, Worrorran languages: \url{https://en.wikipedia.org/wiki/Worrorran_languages}
\item Wikipedia, Wunambal language (treats Miwa as a dialect): \url{https://en.wikipedia.org/wiki/Wunambal_language}
\item Mark Clendon, \emph{Worrorra} (University of Adelaide Press, 2014), a grammar of a related Worrorran language, not of Miwa: \url{https://unglue.it/work/198224/}
\end{itemize}
\section{Religious and Community Resources}
None found.
\section{Digital and Audio Resources}
None found.
\section{Curriculum Development Context}
Any heritage curriculum would have to rely on archival records, held through AIATSIS, and on materials for the closely related Worrorran languages. It should be developed with the relevant Kimberley communities. The Clendon grammar is the only substantial Worrorran reference found.
\end{document}